% ============================================================
%  OPTIQ DSS — Thesis Update Sections
%  To be merged into the main thesis document.
%  New/updated content only; section numbers match the
%  existing structure.
% ============================================================

% -----------------------------------------------------------
%  UPDATED ABSTRACT
% -----------------------------------------------------------
\begin{abstract}
This thesis presents OPTIQ, a web-based Decision Support System (DSS) for
real-time optimisation of the DC4 Butane Debutanizer column at a liquefied
natural gas (LNG) processing facility.  The system integrates an XGBoost
surrogate model, trained on historical distributed control system (DCS) data,
with a Differential Evolution (DE) optimiser to generate energy-efficient
setpoint recommendations subject to purity and operational constraints.

Beyond static recommendation generation, OPTIQ implements a closed-loop
process-tracking workflow.  Upon operator confirmation of a recommendation,
a First-Order Plus Dead Time (FOPDT) simulation engine immediately computes
the predicted process response at three configurable post-apply horizons
(+5, +15, and +30 minutes).  Browser-side timers then fire automatically at
each horizon to compare live DCS readings against the FOPDT predictions.
Deviations exceeding a 5\% threshold trigger a visual alert and prompt
the operator to recompute the recommendation with updated process conditions.

The system is built on a multi-tenant architecture: companies, sites, and
distillation columns are managed through an administrative interface that
allows engineers to configure FOPDT dynamics parameters (gain, time constant,
dead time) and setpoint bounds without modifying source code.  A FastAPI
back-end, React front-end, and PostgreSQL database are deployed on Render's
cloud platform.  Experimental evaluation demonstrates measurable energy
savings while maintaining product purity above specification, together with
a responsive operator experience validated through user testing.
\end{abstract}

% -----------------------------------------------------------
%  SECTION 3.5 — FOPDT Process Dynamics Simulation
% -----------------------------------------------------------
\section{FOPDT Process Dynamics Simulation}
\label{sec:fopdt}

\subsection{Motivation}

Static setpoint recommendations produced by the DE optimiser describe the
expected steady-state improvement but provide no information about the
transient trajectory of the process.  Industrial distillation columns exhibit
significant process lag: changes to reflux ratio, reboiler duty, or feed
temperature do not propagate immediately to the product composition
analyser.  Presenting only a final recommendation therefore leaves the
operator without guidance during the critical minutes immediately following
a setpoint change.

To address this gap, OPTIQ implements a First-Order Plus Dead Time (FOPDT)
simulation layer that computes predicted DCS tag values at each post-apply
horizon and displays them alongside the recommendation.

\subsection{FOPDT Model Formulation}

The FOPDT model characterises the step response of a process variable (PV)
to a change in the manipulated variable (OP) using three parameters:

\begin{itemize}
  \item $K$ — process gain, in PV engineering units per percentage of OP span;
  \item $\tau$ — first-order time constant (seconds), governing how quickly
        the PV approaches its new steady state;
  \item $\theta$ — dead time (seconds), the delay before the PV begins
        to respond.
\end{itemize}

For each output tag $j$ and horizon time $t$ (seconds after apply), the
fractional settling is given by:

\begin{equation}
  \alpha_j(t) =
  \begin{cases}
    0 & t \le \theta_j \\
    1 - \exp\!\left(-\dfrac{t - \theta_j}{\tau_j}\right) & t > \theta_j
  \end{cases}
  \label{eq:fopdt_settling}
\end{equation}

The blended operator output at time $t$ is then:

\begin{equation}
  \mathrm{OP}_j(t) = \mathrm{OP}_j^{\mathrm{cur}}
                   + \alpha_j(t)\,\bigl(\mathrm{OP}_j^{\mathrm{rec}}
                                        - \mathrm{OP}_j^{\mathrm{cur}}\bigr)
  \label{eq:fopdt_blend}
\end{equation}

where $\mathrm{OP}_j^{\mathrm{cur}}$ and $\mathrm{OP}_j^{\mathrm{rec}}$ are
the current and recommended setpoints for tag $j$, respectively.  The
corresponding predicted process variable is:

\begin{equation}
  \mathrm{PV}_j(t) = \mathrm{PV}_j^{\mathrm{nom}}
                   + K_j\,\bigl(\mathrm{OP}_j(t) - 50\bigr)
  \label{eq:fopdt_pv}
\end{equation}

Here $\mathrm{PV}_j^{\mathrm{nom}}$ is the nominal PV value at 50\% OP,
stored as a configurable parameter for each controlled variable.

\subsection{Configured FOPDT Parameters for DC4}

Default FOPDT parameters for the four DC4 manipulated variables are listed
in Table~\ref{tab:fopdt_params}.  These values are representative of typical
debutanizer dynamics and are fully configurable by site engineers through
the Dynamics administration tab without modifying source code.

\begin{table}[h!]
\centering
\caption{Default FOPDT parameters for DC4 Butane Debutanizer}
\label{tab:fopdt_params}
\begin{tabular}{llrrrr}
\hline
\textbf{OP Tag} & \textbf{PV Tag} & \textbf{K} & $\boldsymbol{\tau}$ (s)
  & $\boldsymbol{\theta}$ (s) & \textbf{PV\textsubscript{nom}} \\
\hline
2TIC403.OP & 2TIC403.PV & 0.85 & 1\,200 & 180 & 94.5 \\
2FIC419.OP & 2FIC419.PV & 1.20 & 900  & 120 & 42.0 \\
2LIC409.OP & 2LIC409.PV & 0.60 & 600  &  90 & 55.0 \\
2LIC412.OP & 2LIC412.PV & 0.70 & 750  &  90 & 48.0 \\
\hline
\end{tabular}
\end{table}

\subsection{Simulation Horizons and Output}

The simulation is executed for three configurable horizons: +5, +15, and
+30 minutes (300, 900, and 1\,800 seconds).  At the +30-minute horizon with
$\tau = 1\,200$~s and $\theta = 180$~s, the settling fraction is
approximately 74\%, indicating that the column is still approaching steady
state.  This is expected for slow thermal processes and confirms that a
30-minute follow-up check is essential.

Each horizon produces a \texttt{SimulationStep} object containing the
blended OP, predicted PV, and expected energy and purity values.  These are
returned to the front-end inside the \texttt{ApplyRecommendationOut} response
and rendered as expandable cards in the FOPDT panel.

\subsection{Admin Configurability}

FOPDT parameters and horizon lists are stored in the
\texttt{DistillationColumn.config} JSON field in the database.  The
\texttt{/admin/columns/\{tag\}/fopdt-config} endpoint exposes GET and PUT
operations.  Engineers access these through the Dynamics tab in the
Administration page, where a parameter table allows editing $K$, $\tau$,
$\theta$, and PV\textsubscript{nom} for each OP tag.  Horizon lists are
provided as a comma-separated list of seconds.  Changes take effect
immediately on the next apply operation with no server restart required.

% -----------------------------------------------------------
%  SECTION 3.6 — Closed-Loop Process Tracking
% -----------------------------------------------------------
\section{Closed-Loop Process Tracking}
\label{sec:tracking}

\subsection{Overview}

FOPDT simulation provides predicted trajectories but cannot account for
unmeasured disturbances, feed composition changes, or operating conditions
that differ from the surrogate model's training domain.  Closed-loop process
tracking closes this gap by comparing the FOPDT predictions against actual
DCS readings at each post-apply horizon and informing the operator when the
process has deviated beyond an acceptable tolerance.

\subsection{Timer-Based Trigger Architecture}

Upon receiving the \texttt{ApplyRecommendationOut} response, the React
front-end stores the \texttt{result\_id} and the \texttt{simulation} array
in component state.  It then registers one \texttt{setTimeout} callback per
horizon:

\begin{verbatim}
simulation.forEach((step, idx) => {
  const id = setTimeout(async () => {
    const res = await fopdtConfigApi.checkTracking(
      appliedResultId, currentTagsRef.current
    );
    setTrackingResults(prev => { ... });
  }, step.time_s * 1000);  // fire at exact horizon
});
\end{verbatim}

Live DCS tag values are maintained in a \texttt{useRef} object
(\texttt{currentTagsRef}) that is synchronised with the current WebSocket
broadcast via a \texttt{useEffect} hook.  This pattern avoids the stale
closure problem: the timer callback always reads the tag values that are
current at the moment it fires, not the values captured when the timer
was registered.

\subsection{Deviation Computation}

When a timer fires, the front-end calls
\texttt{POST /recommendations/\{id\}/tracking} with a dictionary of actual
tag values.  The \texttt{ProcessTracker} service:

\begin{enumerate}
  \item Loads the stored \texttt{SimulationStep} list for the result.
  \item Determines the elapsed time since apply.
  \item Interpolates the predicted PV at the elapsed time (selecting the
        nearest simulation step).
  \item Computes, for each PV tag:
  \begin{equation}
    \delta_j = \frac{|\,y_j^{\mathrm{actual}} - y_j^{\mathrm{pred}}\,|}
                    {|\,y_j^{\mathrm{pred}}\,|} \times 100\%
    \label{eq:deviation}
  \end{equation}
  \item Sets \texttt{tracking\_ok = True} if all $\delta_j \le 5\%$.
  \item Sets \texttt{suggest\_reoptimize = True} if any $\delta_j > 5\%$.
\end{enumerate}

The response is a \texttt{TrackingOut} schema containing
\texttt{tracking\_ok}, \texttt{suggest\_reoptimize},
\texttt{worst\_deviation\_pct}, and a \texttt{deviations} dictionary
mapping each PV tag to its predicted value, actual value, and deviation
percentage.

\subsection{User Interface Feedback}

The front-end renders tracking results directly inside each horizon card
(Figure~\ref{fig:tracking_ui}).  When a tracking result arrives:

\begin{itemize}
  \item The card border turns \textbf{green} if \texttt{tracking\_ok},
        or \textbf{red} if a deviation is detected.
  \item The predicted tag value is shown with a strikethrough; the actual
        value is displayed in green (within tolerance) or red (deviation).
  \item A badge summarises the worst deviation:
        ``\checkmark\ DCS tracking · max deviation X\%'' or
        ``\warning\ X\% deviation on [tags]''.
  \item If \texttt{suggest\_reoptimize} is \texttt{True} for any horizon,
        a top-level red banner appears: \textit{``Process deviated from
        expected trajectory. Recompute recommended.''} and the main
        action button changes to ``\warning\ Recompute (conditions changed)''.
\end{itemize}

This design ensures the operator is informed of tracking status with no
manual DCS comparison required, while remaining in full control of the
decision to recompute.

% -----------------------------------------------------------
%  SECTION 4.3 — Updated DSS Architecture
% -----------------------------------------------------------
\section{Updated DSS Architecture}
\label{sec:architecture_updated}

\subsection{Back-End Layer}

The FastAPI back-end is structured around four primary service classes
that collaborate to deliver the optimisation and tracking workflow:

\paragraph{XGBoostSurrogate.} Wraps the serialised XGBoost model and
StandardScaler.  Provides \texttt{predict(readings)} and
\texttt{is\_outlier(readings)} methods used by both the live prediction
endpoint and the DE objective function.

\paragraph{DEOptimiser.} Implements the Differential Evolution algorithm
using \texttt{scipy.optimize.differential\_evolution}.  The objective
function minimises a weighted combination of predicted energy consumption,
purity shortfall, and a move-suppression penalty:
\[
  f(\mathbf{x}) = w_1 E(\mathbf{x}) + w_2 (1 - P(\mathbf{x}))
                + w_3 \cdot \frac{\text{MOVE\_SUPPRESSION}}{n}
                  \sum_{j=1}^{n}\!\left(\frac{\Delta x_j}{s_j}\right)^{\!2}
\]
where $s_j$ is the engineering span of setpoint $j$.  Setpoint bounds are
read from the \texttt{SetpointConfig} entries stored in
\texttt{DistillationColumn.config}.

\paragraph{FopdtSimulator.} Executes Equations~\eqref{eq:fopdt_settling}--\eqref{eq:fopdt_pv}
for each configured horizon and returns a list of \texttt{SimulationStep}
objects.  FOPDT parameters ($K$, $\tau$, $\theta$, PV\textsubscript{nom})
are loaded from \texttt{DistillationColumn.config} at apply time, so
parameter changes by the engineer are immediately reflected without code
changes.

\paragraph{ProcessTracker.} Receives actual DCS tag values from the
front-end at each horizon, compares them against stored
\texttt{SimulationStep} predictions, and returns a \texttt{TrackingResult}
with per-tag deviation percentages and the \texttt{suggest\_reoptimize} flag.

\subsection{Database Schema Extensions}

The following additions to the existing schema support the new workflow:

\begin{itemize}
  \item \textbf{DistillationColumn.config (JSON):}  Extended to store
        \texttt{fopdt\_params} (list of $K$, $\tau$, $\theta$,
        \texttt{pv\_nom} per OP tag) and \texttt{sp\_config} (list of
        setpoint bounds and recommended values).  Horizons are stored under
        the key \texttt{horizons} as a list of integers (seconds).
  \item \textbf{OptimizationResult:}  The \texttt{simulation\_steps}
        field (JSON) stores the serialised list of \texttt{SimulationStep}
        objects computed at apply time, enabling the tracker to retrieve
        predictions without re-running the simulation.
\end{itemize}

\subsection{Front-End Architecture}

The React Dashboard component manages three layers of state that correspond
to the three phases of the workflow:

\begin{enumerate}
  \item \textbf{Recommendation state:}  \texttt{recommendation} (OptimizeOut)
        and \texttt{isStale} (computed from age and tag drift) drive the
        pre-apply panel.
  \item \textbf{Applied state:}  \texttt{applied} (boolean),
        \texttt{appliedResultId} (int), and \texttt{simulation}
        (List[SimulationStep]) are populated when the operator confirms
        the recommendation.
  \item \textbf{Tracking state:}  \texttt{trackingResults}
        (List[TrackingOut | null]) is initialised to \texttt{null} per
        horizon on apply and updated asynchronously as each timer fires.
\end{enumerate}

A \texttt{useRef} hook (\texttt{currentTagsRef}) holds the most recent
DCS tag dictionary and is kept in sync with the live WebSocket feed via a
\texttt{useEffect} dependency.  This pattern ensures that timer callbacks
always read fresh values.

\subsection{Administration Interface}

The Administration page exposes three configuration tabs:

\paragraph{SP Config tab.} Allows engineers to define, for each setpoint
tag, the nominal value, operating bounds (lo/hi), unit, description, and
an optional pre-set recommended value.  Changes update the
\texttt{column.config["sp\_config"]} JSON through
\texttt{PUT /admin/columns/\{tag\}/sp-config}.

\paragraph{Dynamics tab.} Allows engineers to edit FOPDT parameters ($K$,
$\tau$, $\theta$, PV\textsubscript{nom}) for each OP-PV tag pair and
to modify the horizon list.  Changes update
\texttt{column.config["fopdt\_params"]} and \texttt{column.config["horizons"]}
through \texttt{PUT /admin/columns/\{tag\}/fopdt-config}.

\paragraph{Columns tab.} Provides CRUD management for distillation column
records (name, tag, model path, sequence order), accessible to
company-admin and system-admin roles.

\subsection{FOPDT Panel in the Dashboard UI}

After the operator applies a recommendation, the Dashboard renders a
dedicated FOPDT panel below the recommendation summary.  The panel contains
three expandable horizon cards (+5 min, +15 min, +30 min).  Each card shows:

\begin{itemize}
  \item Predicted energy and purity at the horizon, expressed as deltas
        from the current state.
  \item A table of DCS-readable tag values (e.g., 2TIC403 temperature,
        2FIC419 flow rate) with predicted values and a prompt to verify
        on the DCS screen — or, once the tracking timer has fired, the
        actual measured value beside the (struck-through) predicted value,
        colour-coded green or red by deviation magnitude.
  \item A tracking badge summarising the worst deviation and the
        \texttt{suggest\_reoptimize} status for that horizon.
\end{itemize}

The main action button state is determined as follows: during the post-apply
tracking window, the button label is ``\texttt{⚡ Update Recommendation}''
unless at least one horizon's \texttt{TrackingOut.suggest\_reoptimize} is
\texttt{True}, in which case it changes to
``\texttt{⚠ Recompute (conditions changed)}''.  This two-phase logic avoids
false positive warnings during the settling period before any tracking result
has been received.
